\item Analiza la relevancia económica de la minería, la producción de energía y la industria en el mundo.
\item Analiza los efectos de las actividades turísticas en relación con los lugares donde se desarrollan en el mundo.
\item Examina la función del comercio, las redes de comunicaciones y transportes, en la interdependencia económica entre países.
\item Argumenta las desigualdades socioeconómicas en
\item México y el mundo, mediante la interpretación del
\item Índice de Desarrollo Humano (IDH) y el Índice para una Vida Mejor, para explicar sus efectos en la calidad de vida.
\item Comprende qué son las desigualdades socioeconómicas e identifica sus efectos y repercusiones en la vida de las personas.
\item Reconoce cómo las problemáticas sociales y ambientales afectan a la comunidad.
\item Asume responsabilidad como agente de cambio para encontrar soluciones a las problemáticas sociales y ambientales de la comunidad.
\item Comprende qué es la sustentabilidad e identifica prácticas de producción, distribución y consumo sustentables.
\item Argumenta la importancia del consumo responsable, el uso de las tecnologías limpias y los servicios ambientales para contribuir a la sustentabilidad.
\item Argumenta la importancia del consumo responsable, el uso de las tecnologías limpias y los servicios ambientales para contribuir a la sustentabilidad.
\item Comparte alternativas para la protección y recuperación del suelo y colabora de manera organizada y solidaria en acciones comunitarias orientadas a ese fin en la localidad donde vive.